\documentclass[twocolumn,superscriptaddress,aps,pra,longbibliography]{revtex4-2}

\usepackage{amsmath,amssymb,bm}
\usepackage{graphicx}
\usepackage{hyperref}
\usepackage{xcolor}
\usepackage{booktabs}

\newcommand{\bx}{\mathbf{X}}
\newcommand{\br}{\mathbf{r}}
\newcommand{\bv}{\mathbf{v}}
\newcommand{\ket}[1]{|#1\rangle}
\newcommand{\bra}[1]{\langle#1|}
\newcommand{\avg}[1]{\langle #1 \rangle}
\newcommand{\davg}[1]{\langle \Delta #1 \rangle}
\newcommand{\drho}{\hat{\rho}}
\newcommand{\hbar}{\hbar}

\begin{document}

\title{From time-dependent variational Monte Carlo to superfluid hydrodynamics:\\
A frozen-correlation approach for trapped Bose gases}

\author{Xinyuan~Wei}
\affiliation{Department of Physics, University}

\date{\today}

\begin{abstract}
We develop a theoretical and computational framework that connects the
time-dependent variational Monte Carlo (tVMC) method to deterministic
superfluid hydrodynamics for Bose gases in harmonic traps.  Starting from
the full tVMC formalism based on a truncated Jastrow--Feenberg ansatz
$\Psi = e^{\sum_i u_1(\br_i,t) + \sum_{i<j} u_2(r_{ij})}$, we show that
freezing the two-body correlation factor $u_2$ at its ground-state value
reduces the time-dependent variational principle (TDVP) equations to
superfluid hydrodynamics---continuity and Euler equations---with an
effective equation of state $\mu(\rho)$ determined by $u_2$.  This equation
of state is extracted from a one-time variational Monte Carlo calculation
on the uniform system and fit to a cubic polynomial.  The resulting
hydrodynamic equations are solved as a Gross--Pitaevskii equation on a
cylindrical grid with a fourth-order finite-difference Laplacian.  For
$N=8$ microwave-shielded dipolar bosons undergoing a radial trap-frequency
quench ($\omega_\rho: 1\to 2$), the hydrodynamic theory reproduces the
frozen-$u_2$ tVMC breathing dynamics to within 1\% in amplitude and
$<1\%$ in frequency ($\omega_{\rm hydro}=4.19$ vs.\
$\omega_{\rm tVMC}=4.19$), with particle and energy conservation at the
$10^{-3}$ level.  The approach eliminates Monte Carlo sampling from the
real-time evolution loop entirely, reducing computational cost by orders of
magnitude.
\end{abstract}

\maketitle

%======================================================================
\section{Introduction}
\label{sec:intro}
%======================================================================

Understanding the non-equilibrium dynamics of quantum many-body systems
remains a central challenge in ultracold atomic physics.  Recent
experimental advances with microwave-shielded polar
molecules~\cite{Anderegg2023} and magnetic atoms~\cite{Chomaz2023} have
opened new possibilities for studying dipolar Bose--Einstein condensates
(BECs) with tunable interactions, including trap-frequency quenches that
excite collective breathing
modes~\cite{Kinoshita2006,CastinDum1996,Pitaevskii2016}.

From a theoretical standpoint, the time-dependent variational Monte Carlo
(tVMC)
method~\cite{Carleo2012,Carleo2014} provides a powerful approach
to quantum dynamics by combining the expressiveness of correlated
many-body wavefunctions with the time-dependent variational principle
(TDVP)~\cite{Dirac1930,Frenkel1934,McLachlan1964}.
Working with a truncated Jastrow--Feenberg
ansatz~\cite{Jastrow1955,Feenberg1969}, the method captures
both one-body and two-body quantum correlations beyond mean-field theory.
However, the computational cost is substantial: at every time step,
Monte Carlo (MC) sampling is required to estimate the TDVP overlap matrix
and force vector, making long-time simulations expensive.

A key numerical observation motivates the present work: during
trap-frequency quenches that drive large-amplitude breathing oscillations,
the two-body Jastrow factor $u_2(r_{ij})$ remains nearly static while the
one-body field $u_1(\br,t)$ undergoes large changes.
This ``stiffness'' of $u_2$ suggests that short-range two-body correlations
equilibrate rapidly compared to the timescale of density redistribution,
and that it may be possible to freeze $u_2$ at its ground-state value
without significant loss of accuracy.

In this paper, we exploit this observation to derive a complete theoretical
pipeline that eliminates MC from the real-time evolution loop:
\begin{enumerate}
\item We show that freezing $u_2$ and applying the TDVP to the remaining
one-body field $u_1(\br,t)$ yields equations that are \emph{exactly}
equivalent to irrotational superfluid hydrodynamics, with an effective
chemical potential $\mu_{u_2}[\rho]$ determined by the frozen correlations
(Sec.~\ref{sec:hydro_reformulation}).
\item We extract the equation of state (EOS)
$\varepsilon(\rho)$ from a one-time variational MC calculation on the
uniform system with the same frozen $u_2$, and fit it to a polynomial
(Sec.~\ref{sec:eos}).
\item We solve the resulting hydrodynamic equations---cast as a
Gross--Pitaevskii equation (GPE) with the VMC-derived EOS---on a
cylindrical finite-difference grid, and compare quantitatively with
the full frozen-$u_2$ tVMC dynamics for a radial trap quench with $N=8$
particles (Sec.~\ref{sec:results}).
\end{enumerate}
The paper is organized as follows. Section~\ref{sec:model} introduces
the physical model. Sections~\ref{sec:tvmc_formalism}
and~\ref{sec:frozen_u2} present the tVMC formalism and the frozen-$u_2$
approximation with its hydrodynamic reformulation.
Section~\ref{sec:eos} describes the EOS extraction.
Section~\ref{sec:hydro_solver} details the numerical PDE solver.
Section~\ref{sec:results} presents the comparison, and
Sec.~\ref{sec:discussion} discusses the results and outlook.

%======================================================================
\section{Model and Hamiltonian}
\label{sec:model}
%======================================================================

We consider $N$ identical bosons of mass $m$ in a cylindrically symmetric
harmonic trap with time-dependent frequencies,
\begin{equation}
\hat{H}(t) = \sum_{i=1}^{N}\left[
-\frac{\hbar^2}{2m}\nabla_i^2 + V_{\rm tr}(\br_i,t)\right]
+ \sum_{i<j} v(\br_{ij}),
\label{eq:hamiltonian}
\end{equation}
where the trap potential is
\begin{equation}
V_{\rm tr}(\br,t) = \frac{1}{2}m\left[\omega_\rho(t)^2 \rho^2
+ \omega_z^2 z^2\right],
\end{equation}
with $\rho^2 = x^2 + y^2$ the cylindrical radial coordinate.  We work in
natural units $\hbar = m = 1$ and measure lengths in units of the initial
oscillator length $a_{\rm ho} = \sqrt{\hbar/(m\omega_0)}$ with
$\omega_0 = 1$.

\subsection{Interaction potential}

The two-body interaction is modeled as a microwave-shielded dipolar
potential~\cite{Karman2018,Schuster2023} with both long-range dipolar
and short-range shielding contributions:
\begin{equation}
v(\br_{ij}) = V_{\rm dip}(r_{ij},\theta_{ij})
+ V_{\rm shield}(r_{ij},\theta_{ij}),
\label{eq:Vmw}
\end{equation}
where $r_{ij} = |\br_i - \br_j|$ and
$\theta_{ij}$ is the polar angle of the separation vector
$\br_{ij} = \br_i - \br_j$ with respect to the $z$-axis.
The dipolar term has the standard form
\begin{equation}
V_{\rm dip} = \frac{C_3}{r^3}(3\cos^2\theta - 1),
\end{equation}
and the repulsive shielding term is
\begin{equation}
V_{\rm shield} = \frac{C_6}{r^6}
\left(\sin^4\theta + w_{12}\,\sin^2\theta\cos^2\theta\right),
\end{equation}
where $C_3$ and $C_6$ are the dipolar and van der Waals shielding
coefficients, respectively, and $w_{12}$ controls the angular anisotropy
of the shielding.  Both potentials are softened at short range by
$r \to \sqrt{r^2 + \epsilon_r^2}$ with $\epsilon_r = 10^{-3}$ to avoid
numerical singularities.

For the present study we use $N=8$ particles with interaction parameters
$C_3 = 10^{-3}$, $C_6 = 10^{-6}$, and $w_{12} = 2.0$ in oscillator
units, corresponding to a weakly interacting regime where the interaction
energy constitutes approximately 3.5\% of the total ground-state energy.

\subsection{Quench protocol}

The system is initialized in the ground state of an isotropic trap
with $\omega_\rho = \omega_z = 1.0$.  At $t=0$, the radial trap frequency
is instantaneously quenched to $\omega_\rho = 2.0$ while $\omega_z$
remains unchanged.  This excites a radial breathing mode, which we
characterize through the time evolution of the radial moment
$\avg{\rho^2}(t)$ and axial moment $\avg{z^2}(t)$.

%======================================================================
\section{Time-dependent variational Monte Carlo}
\label{sec:tvmc_formalism}
%======================================================================

\subsection{Jastrow--Feenberg wavefunction}

We employ a truncated Jastrow--Feenberg ansatz of the form
\begin{equation}
\Psi(\bx,t) = \exp U(\bx,t), \quad
U(\bx,t) = \sum_{i=1}^{N} u_1(\br_i,t) + \sum_{i<j} u_2(r_{ij},\hat\br_{ij}),
\label{eq:jastrow}
\end{equation}
where $u_1(\br,t)$ is a complex, time-dependent one-body field and
$u_2(r,\hat\br)$ encodes two-body correlations.
The wavefunction is separable in the sense that one-body and two-body
correlations are treated independently, but the full $N$-body probability
distribution $|\Psi|^2$ is non-trivially correlated through the
exponential structure.

\subsubsection{One-body field $u_1$}

In cylindrical coordinates, the one-body field is parametrized with a
separable radial-axial structure:
\begin{equation}
u_1(\rho, z, t) = u_\rho(\rho,t) + u_z(z,t),
\label{eq:u1_separable}
\end{equation}
where each component consists of a harmonic reference plus Gaussian
radial basis function (RBF) corrections:
\begin{align}
u_\rho(\rho,t) &= -\tfrac{1}{2}(\omega_{\rm ref} + a_\rho(t))\,\rho^2
+ \sum_{k=1}^{n_\rho} c_{\rho,k}(t)\, g_0(\rho;\ell_{\rho,k}),
\label{eq:u1_rho}
\\
u_z(z,t) &= -\tfrac{1}{2}(\omega_{\rm ref} + a_z(t))\,z^2
+ \sum_{k=1}^{n_z} c_{z,k}(t)\, g_0(z;\ell_{z,k}).
\label{eq:u1_z}
\end{align}
Here $\omega_{\rm ref}=1.0$ is the reference harmonic frequency,
$a_\rho(t)$ and $a_z(t)$ are complex variational parameters that
adjust the effective trap width, and
$g_0(r;\ell) = \exp(-r^2/(2\ell^2))$ are Gaussian basis functions with
fixed length scales $\ell_{\rho,k}$ and $\ell_{z,k}$ distributed
logarithmically.

In total, the one-body field contains $n_{u_1} = 2 + n_\rho + n_z$
complex variational parameters.  We use $n_\rho = n_z = 6$ with
length scales spanning $\ell \in [0.1, 3.0]$, giving
$n_{u_1} = 14$ parameters.

\subsubsection{Two-body Jastrow $u_2$}

The two-body factor is expanded in partial waves up to $\ell=2$:
\begin{equation}
u_2(r,\cos\theta) = f_0(r) + f_2(r)\,P_2(\cos\theta),
\label{eq:u2_partial}
\end{equation}
where $P_2(x) = (3x^2-1)/2$ is the Legendre polynomial, $\theta$ is the
angle between the interparticle separation and the $z$-axis, and the
radial components are
\begin{align}
f_0(r) &= u_{\rm cusp}(r)
+ \sum_{k=1}^{n_2} c_{0,k}\, g_0(r;\ell_{2,k}),
\label{eq:f0}
\\
f_2(r) &= \sum_{k=1}^{n_2} c_{2,k}\, g_2(r;\ell_{2,k}),
\label{eq:f2}
\end{align}
with $g_2(r;\ell) = r^2\exp(-r^2/(2\ell^2))$ a regular $\ell=2$ basis
function and $n_2=10$ basis functions with length scales
$\ell_2 \in [0.01, 0.7]$.

\subsubsection{Cusp treatment}
\label{sec:cusp}

The short-range cusp function in Eq.~\eqref{eq:f0} is designed to capture
the two-body scattering boundary condition.  We use a Gaussian-damped
form that provides the correct short-range divergence while remaining
computationally tractable:
\begin{equation}
u_{\rm cusp}(r) = -\frac{\kappa}{\sqrt{r^2 + \epsilon_c^2}}\,
\exp\!\left(-\frac{r^2}{2r_c^2}\right),
\label{eq:cusp}
\end{equation}
with cusp strength $\kappa=0.1$, softening parameter
$\epsilon_c = 5\times 10^{-3}$, and cutoff radius $r_c = 0.01$.
The $1/r$ singularity is regularized by $\epsilon_c$ to avoid
numerical overflow, while the Gaussian envelope ensures that the cusp
affects only the short-range pair structure ($r \lesssim r_c$).  The cusp
parameters are fixed (not variational) and determined by the scattering
properties of the two-body potential~\eqref{eq:Vmw}.

The total number of variational parameters is thus
$n_{\rm var} = n_{u_1} + 2n_2 = 14 + 20 = 34$, of which $n_{u_1}=14$
one-body parameters are time-dependent and $2n_2=20$ two-body parameters
are fixed (in the frozen-$u_2$ approximation).

\subsection{Time-dependent variational principle}
\label{sec:tdvp}

The TDVP projects the Schr\"odinger equation onto the variational tangent
space spanned by derivatives of $\ln\Psi$ with respect to the variational
parameters $\{\theta_\alpha\}$.  The resulting equations of motion
are~\cite{Carleo2012}
\begin{equation}
\sum_\beta S_{\alpha\beta}\,\dot\theta_\beta
= -iF_\alpha,
\label{eq:tdvp_general}
\end{equation}
where
\begin{align}
S_{\alpha\beta} &= \avg{\Delta O_\alpha^* \,\Delta O_\beta}_t,
\label{eq:S_def}
\\
F_\alpha &= \avg{\Delta O_\alpha^* \,\Delta E_L}_t,
\label{eq:F_def}
\end{align}
with logarithmic derivatives
$O_\alpha = \partial_{\theta_\alpha}\ln\Psi$,
fluctuations $\Delta A = A - \avg{A}_t$, and the local energy
\begin{equation}
E_L(\bx,t) = \frac{(\hat{H}\Psi)(\bx,t)}{\Psi(\bx,t)}.
\end{equation}

For the Jastrow ansatz~\eqref{eq:jastrow}, the local energy takes the
explicit form
\begin{equation}
E_L = \sum_{i=1}^N\left[
-\tfrac{1}{2}(\nabla_i^2 U + |\nabla_i U|^2) + V_{\rm tr}(\br_i,t)
\right] + \sum_{i<j} v(\br_{ij}),
\label{eq:EL}
\end{equation}
where the gradients of $U = \ln\Psi$ are
\begin{align}
\nabla_i U &= \nabla u_1(\br_i,t) + \sum_{j\neq i}\nabla_i u_2(r_{ij}),
\\
\nabla_i^2 U &= \nabla^2 u_1(\br_i,t) + \sum_{j\neq i}\nabla_i^2 u_2(r_{ij}).
\end{align}

\subsection{Monte Carlo estimation}
\label{sec:mc}

The expectation values in Eqs.~\eqref{eq:S_def}--\eqref{eq:F_def} are
estimated stochastically by sampling configurations $\bx^{(s)}$ from
$|\Psi(\bx,t)|^2$ using the Metropolis algorithm.
We employ single-particle moves: for each walker and each particle~$i$,
a displacement $\Delta\br \sim \text{Uniform}(-\delta/2,\,\delta/2)^3$
is proposed with step size $\delta=0.45$, and accepted with probability
$\min(1,\,|\Psi_{\rm new}|^2/|\Psi_{\rm old}|^2)$.

For each sampled configuration, we compute the observable vector
$O_\alpha^{(s)}$ and local energy $E_L^{(s)}$, from which the overlap
matrix and force vector are constructed as
\begin{align}
S_{\alpha\beta} &= \frac{1}{N_s}\sum_s O_\alpha^{(s)*}O_\beta^{(s)}
- \bar{O}_\alpha^*\bar{O}_\beta,
\\
F_\alpha &= \frac{1}{N_s}\sum_s O_\alpha^{(s)*}E_L^{(s)}
- \bar{O}_\alpha^*\bar{E}_L,
\end{align}
where $\bar{O}_\alpha = N_s^{-1}\sum_s O_\alpha^{(s)}$ and similarly for
$\bar{E}_L$.

The TDVP linear system~\eqref{eq:tdvp_general} is solved via
eigendecomposition of $S$:
\begin{equation}
\dot\theta_\alpha = -i\sum_\beta S^{-1}_{\alpha\beta}F_\beta,
\end{equation}
with eigenvalue truncation at $\lambda_{\min}/\lambda_{\max} = 10^{-4}$
to regularize ill-conditioned directions.  A diagonal shift
$S \to S + \eta\,\mathbb{1}$ with $\eta = 10^{-2}$ provides additional
stability.

\subsection{Time integration}

The equations of motion are integrated using a fourth-order Runge--Kutta
(RK4) scheme.  Each RK4 stage requires independent MC sampling to
evaluate $S$ and $F$ at the intermediate parameter values.
For imaginary-time propagation (ground-state preparation), we use
$\dot\theta = -F/S$ (steepest descent) with $d\tau = 10^{-3}$ for
2000 steps.  For real-time evolution after the quench, we use
$dt = 2\times10^{-4}$ with adaptive time-stepping that limits the
root-mean-square parameter change per step to
$\|\dot\theta\,dt\|_{\rm rms} < 0.01$ and enforces a minimum
effective trap frequency $\omega_{\rm eff} > 0.2$.

The simulation employs $N_w = 65{,}536$ parallel walkers (GPU-accelerated
using JAX~\cite{Bradbury2018}), with $n_{\rm burn} = 30$ equilibration
sweeps followed by $n_{\rm meas} = 10$ measurement sweeps separated by
3~thinning sweeps, per RK4 stage.

%======================================================================
\section{Frozen-$u_2$ approximation and hydrodynamic reformulation}
\label{sec:frozen_u2}
%======================================================================

\subsection{Motivation: stiffness of two-body correlations}

Our numerical tVMC simulations reveal that during the trap quench
$\omega_\rho: 1\to 2$, the two-body Jastrow parameters
$(c_{0,k}, c_{2,k})$ change by less than 1\% of their ground-state
values, while the one-body parameters $(a_\rho, a_z, c_{\rho,k},
c_{z,k})$ undergo order-unity changes.  This reflects the physical
expectation that short-range pair correlations, governed by the
interaction potential on length scales $r \lesssim 1$, equilibrate much
faster than the global density redistribution driven by the trap quench on
scales $r \sim a_{\rm ho}$.

This separation of timescales motivates the \emph{frozen-$u_2$
approximation}: we fix $u_2(r,\hat\br)$ at its ground-state value (real,
time-independent) and evolve only $u_1(\br,t)$ through the TDVP.
The resulting wavefunction
\begin{equation}
\Psi(\bx,t) = \exp\!\left[\sum_i u_1(\br_i,t)
+ \sum_{i<j} u_2(r_{ij})\right]
\label{eq:psi_frozen}
\end{equation}
has the crucial property that the Monte Carlo sampling measure
$|\Psi|^2 \propto \exp[2\sum_i \Re u_1(\br_i,t) + 2\sum_{i<j}u_2(r_{ij})]$
depends only on $\Re u_1$ and the fixed $u_2$, but \emph{not} on
$\Im u_1$.

\subsection{TDVP equations for the active field $u_1$}
\label{sec:tdvp_u1}

When the variational ``parameter'' is the continuous field $u_1(\br,t)$,
the logarithmic derivative becomes a density operator:
\begin{equation}
O(\br;\bx) = \frac{\delta\ln\Psi}{\delta u_1(\br,t)}
= \sum_{i=1}^N \delta^{(3)}(\br - \br_i) \equiv \hat\rho(\br).
\label{eq:O_is_rho}
\end{equation}
The tangent operators are \emph{exactly} the microscopic density operators.

Substituting into the TDVP~\eqref{eq:tdvp_general} yields a continuum
integral equation for $\partial_t u_1$:
\begin{equation}
\boxed{
\int d^3r'\;S(\br,\br';t)\,\partial_t u_1(\br',t)
= -i\,F(\br;t),
}
\label{eq:tdvp_u1}
\end{equation}
where the overlap kernel is the connected density--density correlator
(static structure factor):
\begin{equation}
S(\br,\br';t) = \avg{\Delta\hat\rho(\br)\,\Delta\hat\rho(\br')}_t,
\label{eq:S_rhorho}
\end{equation}
and the force kernel is the density--energy correlator:
\begin{equation}
F(\br;t) = \avg{\Delta\hat\rho(\br)\,\Delta E_L}_t.
\label{eq:F_rhoE}
\end{equation}

The density--density correlator can be decomposed explicitly as
\begin{equation}
S(\br,\br') = \rho(\br)\,\delta^{(3)}(\br-\br')
+ \rho^{(2)}(\br,\br') - \rho(\br)\,\rho(\br'),
\label{eq:S_explicit}
\end{equation}
where $\rho(\br) = \avg{\hat\rho(\br)}_t$ is the one-body density and
$\rho^{(2)}(\br,\br')$ is the (ordered) two-body density.  The frozen
$u_2$ enters through its effect on the sampling measure (which determines
$\rho^{(2)}$) and through the gradient terms in the local
energy~\eqref{eq:EL}.

\emph{Gauge invariance.}---Since $u_1(\br) \to u_1(\br) + c(t)$ only
changes normalization/global phase, $S$ has a corresponding null
direction.  This is handled numerically by imposing a gauge-fixing
constraint, e.g., removing the constant basis function.

\subsection{Reformulation as superfluid hydrodynamics}
\label{sec:hydro_reformulation}

We decompose the complex one-body field into real and imaginary parts:
\begin{equation}
u_1(\br,t) = a(\br,t) + i\theta(\br,t), \qquad a,\theta \in \mathbb{R}.
\label{eq:u1_decomposition}
\end{equation}
Because $u_2$ is real and fixed, the many-body phase is determined
entirely by the one-body phase $\theta$:
\begin{equation}
\arg\Psi(\bx,t) = \sum_{i=1}^N \theta(\br_i,t).
\end{equation}
This implies that the flow on this variational manifold is \emph{purely
one-body and irrotational}, with velocity
\begin{equation}
\bv(\br,t) = \nabla\theta(\br,t)
\label{eq:velocity}
\end{equation}
(in units $\hbar=m=1$).  No backflow terms arise unless time-dependent
two-body phase correlations are allowed.

\subsubsection{Internal energy functional}

Define the \emph{internal energy functional} as the expectation value of
kinetic plus interaction energy in the frozen-$u_2$ Jastrow state at
a given density profile $\rho(\br)$:
\begin{equation}
E_{\rm int}[\rho; u_2] = \avg{\hat{T} + \hat{V}_{\rm int}}_t,
\label{eq:Eint_functional}
\end{equation}
where the expectation is taken in the state~\eqref{eq:psi_frozen} with
$\Re u_1$ chosen to produce the density $\rho$.

If $E_{\rm int}$ can be expressed (or approximated) as a functional of
$\rho$ alone, then the TDVP action reduces to the standard hydrodynamic
form.  The Berry-phase term $\sim \int\rho\,\partial_t\theta$ and the
flow kinetic energy $\sim\frac{1}{2}\int\rho v^2$ combine with the
internal energy to yield the closed pair of PDEs:

\emph{Continuity equation:}
\begin{equation}
\boxed{
\partial_t\rho + \nabla\!\cdot(\rho\,\bv) = 0,
}
\label{eq:continuity}
\end{equation}

\emph{Euler equation:}
\begin{equation}
\boxed{
\partial_t\bv + \nabla\!\left(\tfrac{1}{2}v^2
+ V_{\rm tr}(\br,t) + \mu_{u_2}[\rho]\right) = 0,
}
\label{eq:euler}
\end{equation}
where the effective chemical potential is
\begin{equation}
\mu_{u_2}[\rho](\br,t) = \frac{\delta E_{\rm int}[\rho;u_2]}
{\delta\rho(\br)}.
\label{eq:mu_functional}
\end{equation}

\subsubsection{Relationship to Gross--Pitaevskii theory}

Equations~\eqref{eq:continuity}--\eqref{eq:euler} have exactly the
structure of zero-temperature superfluid hydrodynamics, and differ from
standard Gross--Pitaevskii (GP) theory \emph{only} in the functional form
of $\mu_{u_2}[\rho]$.  In the simplest \emph{local-density approximation}
(LDA), the internal energy takes the form
\begin{equation}
E_{\rm int}^{\rm LDA} = \int d^3r\;\rho(\br)\,
\varepsilon(\rho(\br)),
\label{eq:Eint_LDA}
\end{equation}
where $\varepsilon(\rho)$ is the energy per particle of the uniform system
at density $\rho$ with frozen $u_2$.  The chemical potential then becomes
a local function of density:
\begin{equation}
\mu(\rho) = \frac{d[\rho\,\varepsilon(\rho)]}{d\rho}
= \varepsilon(\rho) + \rho\,\varepsilon'(\rho).
\label{eq:mu_LDA}
\end{equation}

Going beyond LDA, gradient corrections (``quantum pressure'') can be
included.  In the GP framework, these arise from the kinetic energy of
density variations:
\begin{equation}
Q = -\frac{1}{2}\frac{\nabla^2\sqrt{\rho}}{\sqrt{\rho}},
\label{eq:QP}
\end{equation}
which contributes an additional term to the Euler
equation~\eqref{eq:euler}.  As we discuss in Sec.~\ref{sec:hydro_solver},
including this quantum pressure term is essential for accurate dynamics in
the weakly interacting regime.

\subsection{Eliminating Monte Carlo from the time loop}
\label{sec:eliminate_mc}

Even with $u_2$ frozen, the \emph{exact} TDVP
equation~\eqref{eq:tdvp_u1} requires MC sampling at every time step to
evaluate $S(\br,\br';t)$ and $F(\br;t)$, since these depend on the
instantaneous measure $|\Psi(u_1(t), u_2)|^2$.

However, if the assumption of rapid local equilibration holds---meaning
that $u_2$ captures all relevant short-range correlations and that
$E_{\rm int}$ depends on $\rho$ only through the local density---then the
dynamics is fully determined by the EOS $\varepsilon(\rho)$, and no
further MC is needed.  The practical procedure is:
\begin{enumerate}
\item Perform a \emph{one-time} VMC calculation on the uniform system
(periodic boundary conditions, frozen $u_2$) at a grid of densities to
extract $\varepsilon(\rho)$.
\item Fit $\varepsilon(\rho)$ to a polynomial to obtain a smooth, analytic
$\mu(\rho)$.
\item Solve the hydrodynamic equations~\eqref{eq:continuity}--\eqref{eq:euler}
(or their GP reformulation) deterministically on a spatial grid.
\end{enumerate}
This replaces the $O(N_{\rm steps} \times N_{\rm MC})$ cost of tVMC with
$O(N_{\rm MC})$ (one-time EOS) + $O(N_{\rm steps} \times N_{\rm grid})$
(PDE evolution), yielding orders-of-magnitude speedup for the real-time
dynamics.

%======================================================================
\section{Equation of state from VMC}
\label{sec:eos}
%======================================================================

\subsection{Uniform-system VMC}

To extract the EOS, we perform VMC calculations for a uniform system of
$N=8$ particles in a cubic box with periodic boundary conditions, at 20
density values logarithmically spaced over $\rho \in [0.05, 3.0]$.  The
two-body Jastrow $u_2(r,\hat\br)$ is frozen at the ground-state values
obtained from the trapped-system tVMC calculation, while $u_1$ is taken as
constant (since the density is uniform).  The one-body field has no
variational freedom in the uniform case; it merely sets the overall
normalization.

At each density, we use $N_w = 8{,}192$ walkers with 200 equilibration
sweeps and 500 measurement sweeps (thinning by 5) to estimate the energy
per particle $\varepsilon(\rho)$.  The kinetic energy is computed using the
Hermitian estimator $T = \frac{1}{2}\avg{|\nabla U|^2 + \nabla^2 U}$ to
ensure correct results.  Error bars are obtained from a blocking analysis
over the 500 measurement batches.

\subsection{Polynomial fitting}
\label{sec:eos_fit}

We fit the VMC data to a cubic polynomial in density:
\begin{equation}
\varepsilon(\rho) = a_1\rho + a_2\rho^2 + a_3\rho^3,
\label{eq:eps_poly3}
\end{equation}
using weighted least-squares with the VMC error bars as weights.
Model selection between quadratic ($a_3=0$) and cubic forms is
performed using the Akaike information criterion (AIC), which favors
the cubic with $\Delta\text{AIC} > 2$.

The resulting fit parameters are
\begin{equation}
\begin{aligned}
a_1 &= 0.1017 \pm 0.0035, \\
a_2 &= 0.0232 \pm 0.0062, \\
a_3 &= -0.0067 \pm 0.0019,
\end{aligned}
\label{eq:eos_coeffs}
\end{equation}
with reduced chi-squared $\chi^2_\nu = 0.41$, indicating a good fit.
The chemical potential follows from Eq.~\eqref{eq:mu_LDA}:
\begin{equation}
\mu(\rho) = 2a_1\rho + 3a_2\rho^2 + 4a_3\rho^3.
\label{eq:mu_poly3}
\end{equation}

The negative cubic coefficient $a_3 < 0$ introduces a stability limit:
$d\mu/d\rho = 0$ at
$\rho_{\rm stab} \approx 2.67$, above which the compressibility becomes
negative.  For the densities encountered in our simulations
($\rho_{\max} \approx 5$ at maximum compression), we regularize the EOS by
linearly extrapolating $\mu(\rho)$ for $\rho > 0.8\,\rho_{\rm stab}$ with
a minimum slope of $2a_1$, ensuring thermodynamic stability at all
densities.

%======================================================================
\section{Hydrodynamic PDE solver}
\label{sec:hydro_solver}
%======================================================================

\subsection{Gross--Pitaevskii formulation}

Rather than solving the hydrodynamic
equations~\eqref{eq:continuity}--\eqref{eq:euler} directly in the Madelung
variables $(\rho,\theta)$, we reformulate them as a Gross--Pitaevskii
equation (GPE) for the complex field
$\psi(\br,t) = \sqrt{\rho}\,e^{i\theta}$:
\begin{equation}
\boxed{
i\,\partial_t\psi = \left[-\tfrac{1}{2}\nabla^2
+ V_{\rm tr}(\br,t) + \mu(|\psi|^2)\right]\psi,
}
\label{eq:gpe}
\end{equation}
where $\mu(\rho)$ is the VMC-derived chemical potential~\eqref{eq:mu_poly3}
and the $-\frac{1}{2}\nabla^2$ term provides both the quantum pressure
and the kinetic energy of the flow.

This GPE formulation offers several numerical advantages over the Madelung
form: (i)~it naturally includes the quantum pressure without explicit
computation of $\nabla^2\sqrt\rho/\sqrt\rho$, which is singular where
$\rho\to 0$; (ii)~it avoids the velocity-field singularity at vacuum
boundaries; and (iii)~it provides a single complex field equation instead
of two coupled real equations.

\emph{Note on the quantum pressure.}---In the strict frozen-$u_2$ TDVP
derivation (Sec.~\ref{sec:hydro_reformulation}), the quantum pressure
arises from the density dependence of $S(\br,\br')$ beyond the LDA.  For
the weakly interacting $N=8$ system considered here, the quantum pressure
is numerically important: without it, the dynamics is qualitatively
different.  We therefore include the full Laplacian in
Eq.~\eqref{eq:gpe}.

\subsection{Cylindrical coordinates}

Exploiting the azimuthal symmetry of the problem, we reduce the 3D GPE
to a 2D equation in cylindrical coordinates $(r,z)$, where the Laplacian
is
\begin{equation}
\nabla^2\psi = \frac{\partial^2\psi}{\partial r^2}
+ \frac{1}{r}\frac{\partial\psi}{\partial r}
+ \frac{\partial^2\psi}{\partial z^2}.
\label{eq:lap_cyl}
\end{equation}

The computational domain is $r \in [0, R_{\max}]$,
$z \in [-Z_{\max}, Z_{\max}]$ with $R_{\max} = Z_{\max} = 5.0$.

\emph{Boundary conditions:}
At $r=0$, azimuthal symmetry requires $\partial_r\psi|_{r=0} = 0$.
We implement this via a ghost point $\psi(-\Delta r) = \psi(\Delta r)$
(even extension).  The $(1/r)\partial_r\psi$ term is evaluated at $r=0$
using L'H\^opital's rule:
$\lim_{r\to 0}(1/r)\partial_r\psi = \partial_r^2\psi|_{r=0}$.
At the outer boundaries ($r = R_{\max}$, $z = \pm Z_{\max}$), we impose
Dirichlet conditions $\psi = 0$, which is consistent with the wavefunction
vanishing far from the trap center.

\subsection{Fourth-order spatial discretization}

We use a uniform grid with $N_r \times N_z$ points.  Interior spatial
derivatives are computed using fourth-order central finite differences:
\begin{align}
\frac{\partial^2 f}{\partial x^2}\bigg|_i
&= \frac{-f_{i-2} + 16f_{i-1} - 30f_i + 16f_{i+1} - f_{i+2}}{12\,\Delta x^2}
\nonumber\\
&\qquad + O(\Delta x^4),
\label{eq:fd4_d2}
\\[4pt]
\frac{\partial f}{\partial x}\bigg|_i
&= \frac{-f_{i+2} + 8f_{i+1} - 8f_{i-1} + f_{i-2}}{12\,\Delta x}
\nonumber\\
&\qquad + O(\Delta x^4).
\label{eq:fd4_d1}
\end{align}
Near boundaries, ghost-point values are determined by symmetry ($r=0$) or
the Dirichlet condition ($\psi=0$ outside the domain).  The fourth-order
scheme significantly improves the accuracy of the breathing frequency
compared to second-order differences (Sec.~\ref{sec:convergence}).

\subsection{Time integration}

The GPE~\eqref{eq:gpe} is integrated in time using an explicit RK4
scheme.  Writing $\psi_{n+1} = \psi_n + (dt/6)(k_1 + 2k_2 + 2k_3 + k_4)$
with
\begin{equation}
k_j = -i\!\left[-\tfrac{1}{2}\nabla^2 + V + \mu(|\cdots|^2)\right]
(\psi_n + \text{RK increment}),
\end{equation}
the time step is determined by the dispersive CFL condition
\begin{equation}
\Delta t \leq C_{\rm CFL}\,\frac{(\min\Delta x)^2}{\pi},
\label{eq:cfl}
\end{equation}
with $C_{\rm CFL} = 0.3$.  This ensures stability for the $-\frac{i}{2}\nabla^2$
dispersive term.  An additional potential constraint
$\Delta t \leq 2\,C_{\rm CFL}/|\mu_{\max}|$ limits the phase
accumulation per step.

\subsection{Initial conditions}

We implement two options for the initial wavefunction $\psi(r,z,0)$:

\emph{(i)~Gaussian IC:}
A Gaussian density profile matched to the tVMC ground-state moments
$\avg{\rho^2}_0$ and $\avg{z^2}_0$:
\begin{equation}
\rho_0(r,z) \propto \exp\!\left(-\frac{r^2}{2\sigma_r^2}
- \frac{z^2}{2\sigma_z^2}\right),
\end{equation}
with $\sigma_r^2 = \avg{\rho^2}_0/2$ and $\sigma_z^2 = \avg{z^2}_0$,
normalized to $\int\rho_0\,2\pi r\,dr\,dz = N$.

\emph{(ii)~GPE ground state:}
The self-consistent ground state of Eq.~\eqref{eq:gpe} in the
pre-quench trap, obtained by imaginary-time propagation:
\begin{equation}
\psi \to \psi - d\tau\,\hat{H}\psi, \qquad
\psi \to \psi\sqrt{N/\|\psi\|^2},
\end{equation}
iterated until the chemical potential $\mu = \avg{\psi^*\hat{H}\psi}/N$
converges to $\delta\mu < 10^{-8}$.  Starting from a Gaussian guess,
convergence is typically achieved in $\sim 30{,}000$ iterations with
$d\tau = 0.3\,\Delta x^2/\pi$.

The GPE ground state is the more physical choice, as it is self-consistent
with the equations of motion and produces the correct initial energy.
For our system, it yields $\avg{\rho^2}_0 = 1.058$, $\avg{z^2}_0 = 0.529$,
and total energy $E_0 = 12.44$, in excellent agreement with the tVMC
ground-state values (Table~\ref{tab:comparison}).

\subsection{Observables}

At each time step (or every 10 steps for efficiency), we compute the
following observables using cylindrical trapezoidal integration
$\int f\,2\pi r\,dr\,dz$:
\begin{align}
N(t) &= \int |\psi|^2\,2\pi r\,dr\,dz, \label{eq:obs_N}\\
\avg{\rho^2}(t) &= \frac{1}{N}\int r^2\,|\psi|^2\,2\pi r\,dr\,dz,\\
\avg{z^2}(t) &= \frac{1}{N}\int z^2\,|\psi|^2\,2\pi r\,dr\,dz,\\
E(t) &= \int\!\left[-\tfrac{1}{2}\psi^*\nabla^2\psi
+ V|\psi|^2 + |\psi|^2\varepsilon(|\psi|^2)\right]
2\pi r\,dr\,dz.
\label{eq:obs_E}
\end{align}
The kinetic energy in $E(t)$ is computed as
$-\frac{1}{2}\Re\int\psi^*\nabla^2\psi$ using the same discrete Laplacian
as the dynamics, ensuring consistent energy conservation.

\subsection{Scaling ansatz}

As a further check, we compare with the scaling ansatz for breathing modes
in a harmonic trap~\cite{CastinDum1996,Kagan1996}.  Assuming a
self-similar density evolution
$\rho(r,z,t) = (\lambda_r^2\lambda_z)^{-1}\rho_0(r/\lambda_r,z/\lambda_z)$,
the scaling parameters satisfy
\begin{align}
\ddot\lambda_r &= -\omega_\rho^2\lambda_r
+ \frac{C_r}{\lambda_r^{2\gamma+1}\lambda_z^\gamma},
\label{eq:scaling_r}\\
\ddot\lambda_z &= -\omega_z^2\lambda_z
+ \frac{C_z}{\lambda_r^{2\gamma}\lambda_z^{\gamma+1}},
\label{eq:scaling_z}
\end{align}
where $C_r = \omega_{\rho,0}^2$ and $C_z = \omega_{z,0}^2$ for the
non-interacting ($\gamma=1$) limit.

%======================================================================
\section{Results}
\label{sec:results}
%======================================================================

\subsection{Ground-state properties}

Table~\ref{tab:comparison} compares the ground-state properties from
three approaches: full tVMC (all parameters evolving), frozen-$u_2$ tVMC
(20 two-body parameters fixed), and the GPE with the VMC-derived EOS.

\begin{table}[t]
\caption{Ground-state properties for $N=8$ in an isotropic trap
($\omega_\rho = \omega_z = 1$).  ``Full tVMC'' evolves all 30
variational parameters; ``Frozen-$u_2$ tVMC'' fixes 20 two-body parameters
and evolves only 14 one-body parameters; ``GPE'' is the hydrodynamic
solver with the VMC-derived EOS (imaginary-time ground state).}
\label{tab:comparison}
\begin{ruledtabular}
\begin{tabular}{lccc}
& Full tVMC & Frozen-$u_2$ & GPE \\
\hline
$E_{\rm gs}$ & 12.433 & 12.423 & 12.442 \\
$\avg{\rho^2}_0$ & 1.049 & 1.051 & 1.058 \\
$\avg{z^2}_0$ & 0.525 & 0.525 & 0.529 \\
Aspect $\sqrt{\avg{\rho^2}/(2\avg{z^2})}$ & 0.999 & 1.000 & 1.001 \\
\end{tabular}
\end{ruledtabular}
\end{table}

The frozen-$u_2$ tVMC matches the full tVMC to within 0.1\% in energy
and $<$0.3\% in spatial moments, confirming that $u_2$-freezing does not
significantly affect the ground state.  The GPE ground state agrees with
the frozen-$u_2$ tVMC to 0.15\% in energy and $<$0.7\% in moments.  The
small discrepancy reflects the combined effect of the LDA approximation,
the polynomial EOS truncation, and the finite-grid discretization.

\subsection{Quench dynamics: breathing oscillations}

Figure~\ref{fig:breathing} shows the time evolution of $\avg{\rho^2}(t)$
and $\avg{z^2}(t)$ after the quench $\omega_\rho: 1\to 2$ at $t=0$.

\begin{figure}[t]
\includegraphics[width=\columnwidth]{hydro_comparison.png}
\caption{Breathing dynamics after the radial trap quench
$\omega_\rho: 1\to 2$ for $N=8$ particles.
\textbf{Top:} Radial moment $\avg{\rho^2}(t)$; red: frozen-$u_2$ tVMC,
blue: GPE with VMC EOS (ground-state IC), green dashed: GPE with
Gaussian IC, gray dotted: scaling ansatz.
\textbf{Middle:} Axial moment $\avg{z^2}(t)$.
\textbf{Bottom:} Conservation diagnostics---relative particle-number and
energy errors ($128\times 256$ grid, fourth-order Laplacian).}
\label{fig:breathing}
\end{figure}

The radial moment oscillates between $\avg{\rho^2}_{\min} \approx 0.26$
and $\avg{\rho^2}_{\max} \approx 1.06$, representing a fourfold
compression in the radial direction.  The axial moment shows a weaker
sympathetic response, oscillating between
$\avg{z^2}_{\min} \approx 0.53$ and $\avg{z^2}_{\max} \approx 0.61$.

The GPE dynamics (both initial conditions) tracks the frozen-$u_2$
tVMC reference closely throughout the simulation:
\begin{itemize}
\item The oscillation amplitude in $\avg{\rho^2}$ agrees to within 1\%:
GPE range $[0.263, 1.058]$ vs.\ tVMC $[0.264, 1.055]$.
\item The $\avg{z^2}$ evolution matches to within 1\%: GPE range
$[0.529, 0.607]$ vs.\ tVMC $[0.522, 0.613]$.
\item The energy jump at the quench is
$\Delta E = \frac{1}{2}(\omega_{\rho,f}^2 - \omega_{\rho,i}^2)
\avg{\rho^2}_0\,N = 12.70$ (GPE ground-state IC) vs.\ the expected
value $12.70$, confirming exact virial consistency.
\end{itemize}

\subsection{Breathing frequency}

The dominant oscillation frequency is extracted from the Fourier transform
of $\avg{\rho^2}(t)$.  Both the frozen-$u_2$ tVMC and the GPE yield
\begin{equation}
\omega_{\rm br} = 4.19 \approx 2\omega_\rho,
\label{eq:omega_br}
\end{equation}
consistent with the $2\omega$ breathing mode expected for a 2D
harmonic confinement.

The scaling ansatz~\eqref{eq:scaling_r}--\eqref{eq:scaling_z} also
predicts a breathing frequency close to $2\omega_\rho$ but with a
larger oscillation amplitude (because it neglects quantum pressure and
anharmonic effects), as visible in Fig.~\ref{fig:breathing}.

\emph{Role of quantum pressure.}---The quantum pressure
term~\eqref{eq:QP} is essential: without it, the GPE reduces to a
potential-only evolution that does not properly couple density and velocity,
and the dynamics is qualitatively incorrect for this weakly interacting
($E_{\rm int}/E_{\rm tot} \approx 3.5\%$) system.

\subsection{Conservation and convergence}
\label{sec:convergence}

The GPE conserves particle number and energy exactly in the continuum
limit.  On a discrete grid, conservation errors arise from spatial
discretization.  Table~\ref{tab:conservation} summarizes the conservation
diagnostics for different grid resolutions and Laplacian orders.

\begin{table}[t]
\caption{Conservation diagnostics for the GPE simulation over
$t \in [0, 3.0]$.  $\delta N/N$ and $\delta E/E$ are the maximum relative
deviations of particle number and energy from their initial values.}
\label{tab:conservation}
\begin{ruledtabular}
\begin{tabular}{lccc}
Grid & $\nabla^2$ order & $\max|\delta N/N|$ & $\max|\delta E/E|$ \\
\hline
$128\times256$ & 2nd & $1.09\times10^{-3}$ & $1.20\times10^{-3}$ \\
$128\times256$ & 4th & $7.30\times10^{-4}$ & $1.29\times10^{-3}$ \\
$160\times320$ & 4th & $4.65\times10^{-4}$ & $8.22\times10^{-4}$ \\
\end{tabular}
\end{ruledtabular}
\end{table}

Several observations are noteworthy:
\begin{enumerate}
\item The conservation errors are \emph{purely spatial}:
halving the CFL factor from 0.3 to 0.15 (doubling the number of time
steps) produces identical conservation to machine precision, confirming
that the RK4 temporal error is negligible.
\item The errors are \emph{oscillatory}, peaking during maximum compression
(when the density gradients are steepest) and nearly vanishing when the
cloud re-expands to its initial size.  At $t=1.5$ (full expansion),
$|\delta N/N| = 3.2\times10^{-6}$ and
$|\delta E/E| = 2.8\times10^{-6}$.
\item The fourth-order Laplacian significantly improves the breathing
frequency accuracy: the second-order scheme gives $\omega_{\rm br} = 3.99$,
while the fourth-order scheme yields $\omega_{\rm br} = 4.19$, matching
the tVMC value.
\end{enumerate}

\subsection{Summary of the comparison}

Table~\ref{tab:dynamics} provides a comprehensive comparison of the
dynamical observables.

\begin{table}[t]
\caption{Comparison of quench dynamics for $\omega_\rho: 1\to 2$, $N=8$.
GPE results are for the ground-state IC on a $128\times256$ grid with
fourth-order Laplacian.}
\label{tab:dynamics}
\begin{ruledtabular}
\begin{tabular}{lcc}
& Frozen-$u_2$ tVMC & GPE \\
\hline
$E_{\rm gs}$ & 12.423 & 12.442 \\
$\avg{\rho^2}_{\min}$ & 0.264 & 0.263 \\
$\avg{\rho^2}_{\max}$ & 1.055 & 1.058 \\
$\avg{z^2}_{\min}$ & 0.522 & 0.529 \\
$\avg{z^2}_{\max}$ & 0.613 & 0.607 \\
$\omega_{\rm br}$ & 4.19 & 4.19 \\
$\Delta E/N$ at quench & 1.59 & 1.59 \\
\end{tabular}
\end{ruledtabular}
\end{table}

The agreement is excellent across all observables, demonstrating that
the hydrodynamic theory with a VMC-derived EOS faithfully reproduces
the frozen-$u_2$ tVMC dynamics for this system.

%======================================================================
\section{Discussion and outlook}
\label{sec:discussion}
%======================================================================

\subsection{Validity of the frozen-$u_2$ approximation}

The success of the frozen-$u_2$ framework rests on two conditions:
(i)~the two-body correlations must be ``stiff'' (slowly varying compared
to one-body dynamics), and (ii)~the EOS must adequately capture the
density-dependent energetics.

For the weakly interacting $N=8$ system studied here, both conditions are
well satisfied.  The interaction energy is only $\sim$3.5\% of the total
energy, so the dominant dynamics is kinetic (quantum pressure) and
potential (trap).  The frozen-$u_2$ tVMC agrees closely with the full
tVMC in terms of oscillation amplitudes ($\avg{\rho^2}$ range
$[0.264, 1.055]$ vs.\ $[0.264, 1.060]$), though the frequency spectra
differ in detail due to the additional degrees of freedom in the full
calculation.

\subsection{Role of quantum pressure}

A key finding is that the quantum pressure (Bohm potential) is essential
for accurate dynamics in the present regime.  Without it, the Madelung
hydrodynamic equations are either unstable at cloud boundaries (where
$\rho\to 0$ creates singularities) or predict incorrect breathing
amplitudes.  The GPE formulation naturally incorporates quantum pressure
through the $-\frac{1}{2}\nabla^2\psi$ term, avoiding both problems.

This is expected physically: for $N=8$ with weak interactions, the
quantum pressure provides the dominant repulsive mechanism that
prevents over-compression.  In the Thomas--Fermi limit of large $N$
and strong interactions, the EOS contribution $\mu(\rho)$ would dominate,
and the quantum pressure could be neglected.

\subsection{Computational cost comparison}

The frozen-$u_2$ tVMC simulation with 65,536 walkers, 30 burn-in sweeps,
and 10 measurement sweeps per RK4 stage requires $\sim$6 hours on a
modern GPU for $t_{\max}=3.0$.  In contrast, the GPE solver completes
in $\sim$4 minutes (including $\sim$3 minutes for imaginary-time
ground-state preparation and $\sim$1 minute for real-time evolution) on a
single CPU core.  The one-time EOS computation adds $\sim$30 minutes.
The total cost reduction is approximately two orders of magnitude.

\subsection{Outlook: larger $N$ and stronger interactions}

The hydrodynamic approach is expected to become \emph{more} accurate at
larger $N$ and stronger interactions, where:
\begin{itemize}
\item The Thomas--Fermi limit suppresses quantum pressure effects, making
the LDA more accurate.
\item The EOS captures a larger fraction of the total energetics.
\item The density profile is smoother (less steep gradients), improving
the finite-difference accuracy.
\end{itemize}
A natural next step is to validate the framework for $N=16$, 32, and
beyond, where the full tVMC becomes increasingly expensive but the
hydrodynamic approach remains fast.  This would establish the range of
applicability and identify the crossover from quantum-pressure-dominated
to EOS-dominated dynamics.

Another direction is to go beyond the LDA by incorporating gradient
corrections to the EOS, either from analytical theory or by fitting
inhomogeneous VMC data.  This could improve accuracy at small $N$ where
the LDA assumption breaks down.

%======================================================================
\section{Conclusion}
\label{sec:conclusion}
%======================================================================

We have established a complete theoretical and computational pipeline
connecting time-dependent variational Monte Carlo to deterministic
superfluid hydrodynamics for trapped Bose gases.  The key steps are:

\begin{enumerate}
\item \emph{Theory:} Freezing the two-body Jastrow factor $u_2$ and
applying the TDVP to the one-body field $u_1(\br,t)$ yields equations
that are exactly equivalent to irrotational superfluid hydrodynamics,
with an effective chemical potential determined by the frozen correlations.

\item \emph{Equation of state:} A one-time VMC calculation on the uniform
system extracts the density-dependent energy per particle
$\varepsilon(\rho) = 0.102\rho + 0.023\rho^2 - 0.007\rho^3$, from
which the chemical potential $\mu(\rho)$ follows analytically.

\item \emph{Hydrodynamic solver:} The GPE with the VMC-derived EOS,
solved on a cylindrical grid with a fourth-order Laplacian and RK4
time integration, reproduces the frozen-$u_2$ tVMC breathing dynamics
to within 1\% in amplitude and $<$1\% in frequency for $N=8$
microwave-shielded dipolar bosons undergoing a radial trap quench.
\end{enumerate}

The approach reduces computational cost by two orders of magnitude by
eliminating Monte Carlo sampling from the real-time evolution loop.
It provides a practical framework for studying collective dynamics in
trapped quantum gases, with clear pathways to extension for larger
systems and stronger interactions.

\begin{acknowledgments}
Computations were performed using JAX~\cite{Bradbury2018} for
GPU-accelerated VMC and NumPy/SciPy for the hydrodynamic solver.
\end{acknowledgments}

\bibliography{references}

\end{document}
